% complexification.tex — Section 2: Part A — Complexification from C*-Observer
% ASSERT_CONVENTION: natural_units=dimensionless, jordan_product=(1/2)(ab+ba), clifford_convention=euclidean_positive, octonion_basis=fano_e1e2=e4, complex_structure=u_equals_e7, spin_representation=S10plus_boyle

%----------------------------------------------------------------------
\section{Part A: Complexification from C*-Observer}\label{sec:complexification}
%----------------------------------------------------------------------

The first half of the derivation proves that the observer's algebraic
nature---specifically, the fact that self-modeling forces a C*-algebra
structure---produces complexification of the Peirce half-space of the
exceptional Jordan algebra, upgrading the symmetry from
$\Spin{9}$ to $\Spin{10}$ and identifying the 16-dimensional
spinor representation that will carry one generation of fermions.

%----------------------------------------------------------------------
\subsection{The exceptional Jordan algebra \texorpdfstring{$\hthree$}{h3(O)}}\label{sec:h3O}
%----------------------------------------------------------------------

The exceptional Jordan algebra $\hthree$ (the Albert algebra)
consists of $3\times 3$ Hermitian matrices over the octonions $\mathbb{O}$:
\begin{equation}\label{eq:h3O-element}
  X \;=\;
  \begin{pmatrix}
    \alpha & \bar{x}_3 & x_2 \\[3pt]
    x_3 & \beta & \bar{x}_1 \\[3pt]
    \bar{x}_2 & x_1 & \gamma
  \end{pmatrix},
  \qquad
  \alpha,\beta,\gamma \in \mathbb{R},
  \quad
  x_1,x_2,x_3 \in \mathbb{O},
\end{equation}
with the Jordan product
\begin{equation}\label{eq:jordan-product}
  A \circ B \;=\; \tfrac{1}{2}(AB + BA),
\end{equation}
where $AB$ denotes ordinary matrix multiplication with octonion entries.
The product is well-defined even though $\mathbb{O}$ is non-associative,
by the classical result of Albert~\cite{Albert1934}; the exceptional
Jordan algebra was classified by Jordan, von~Neumann, and
Wigner~\cite{JvNW1934}.

\begin{remark}[Dimension]
The three real diagonal entries contribute $3$; the three independent
octonion off-diagonal entries contribute $3\times 8 = 24$.  Hence
$\dim(\hthree) = 3 + 24 = 27$.
\end{remark}

The identification of $\hthree$ as the algebraic basin in which a
self-modeling observer lives is packaged here as a Definition plus a
Theorem.  The Definition states the structural conditions any
candidate basin must satisfy; the Theorem is that, up to isomorphism,
$\hthree$ is the unique candidate.  An interpretive remark
relating the conditions to the Level~IV realism premise is recorded
afterwards (Remark~\ref{rem:basin-scope}); it is upstream motivation,
not part of the proof.

\begin{definition}[Observer-grade JA basin]\label{def:basin}
A simple finite-dimensional formally real Jordan algebra $A$ is an
\emph{observer-grade JA basin} if all three conditions hold:
\begin{enumerate}[label=(\roman*)]
  \item $A$ admits a Peirce decomposition at a rank-1 idempotent $E$
    whose Peirce 0-space (spectator complement) $V_0(E)$ is
    isomorphic to $h_2(\mathbb{K})$ for some composition algebra
    $\mathbb{K}\in\{\mathbb{R},\mathbb{C},\mathbb{H},\mathbb{O}\}$;
  \item $\mathbb{K}$ admits a complex structure: there exists
    $u\in\mathbb{K}$ with $u^2=-1$, so that
    $\mathbb{C}_u=\mathbb{R}\oplus\mathbb{R}u\subseteq\mathbb{K}$
    is a complex subfield and $h_2(\mathbb{C}_u)\cong
    M_2(\mathbb{C})^{sa}$ sits inside $h_2(\mathbb{K})$ as a
    (Jordan-)subalgebra;
  \item $A$ is not a non-trivial Jordan-tensor factor in the
    Barnum--Graydon--Wilce~\cite{BGW2020} symmetric monoidal
    category of formally real JC-algebras (special FRJAs).
\end{enumerate}
\end{definition}

\noindent
Condition~(i) restricts attention to rank-3 algebras over a
composition algebra (it fails for spin factors of rank~2 and for
$h_n(\mathbb{K})$ with $n\geq 4$, where $V_0$ has the wrong shape).
Condition~(ii) excludes the case $\mathbb{K}=\mathbb{R}$, where no
complex sub-line exists and the C*-bottleneck of
Lemma~\ref{lem:bottleneck} below is not defined; for
$\mathbb{K}\in\{\mathbb{C},\mathbb{H},\mathbb{O}\}$ a complex
structure exists ($u=i$, $u\in\{i,j,k\}$, or $u\in S^6$,
respectively).  Condition~(iii) is the categorical
non-composability statement: it excludes the special simple FRJAs,
which are the objects of the BGW category, leaving only $\hthree$.

\begin{theorem}[Restricted Basin]\label{thm:basin}
Up to isomorphism, the unique observer-grade JA basin is $\hthree$.
\end{theorem}

\begin{proof}
By the JvNW classification~\cite{JvNW1934}, every simple finite-
dimensional formally real Jordan algebra is one of $h_n(\mathbb{R})$,
$h_n(\mathbb{C})$, $h_n(\mathbb{H})$ for $n\geq 3$, the spin factors
$V_k$ for $k\geq 2$, or $\hthree$.  We check each family against
Definition~\ref{def:basin}.
\begin{enumerate}
  \item \textbf{Spin factors $V_k$ ($k\geq 2$) fail (i).}
    A spin factor has rank~$2$; for any rank-1 idempotent the
    Peirce 0-space is the orthogonal rank-1 idempotent's span,
    which is one-dimensional and not of the form $h_2(\mathbb{K})$
    for any composition algebra.
  \item \textbf{$h_n(\mathbb{K})$ for $n\geq 4$,
    $\mathbb{K}\in\{\mathbb{R},\mathbb{C},\mathbb{H}\}$, fail (i).}
    For $n\geq 4$ a rank-1 idempotent has Peirce 0-space
    $V_0(E)\cong h_{n-1}(\mathbb{K})$, which is not of the form
    $h_2(\mathbb{K})$.
  \item \textbf{$h_3(\mathbb{R})$ fails (ii).}  $\mathbb{K}=\mathbb{R}$
    contains no element $u$ with $u^2=-1$, so no complex sub-line
    $\mathbb{C}_u\subseteq\mathbb{R}$ exists; condition~(ii) is
    vacuously violated.
  \item \textbf{$h_3(\mathbb{C})$ and $h_3(\mathbb{H})$ fail (iii).}
    Both are special simple FRJAs, equivalently formally real
    JC-algebras: each embeds as the self-adjoint part of an
    associative matrix algebra ($M_3(\mathbb{C})$ and $M_3(\mathbb{H})$
    respectively) with Jordan product $a\circ b=\tfrac12(ab+ba)$.
    By BGW Theorem~4.15~\cite{BGW2020}, simple special EJAs $A,B$
    admit a composite $AB$ that is an ideal in the Hanche-Olsen
    universal tensor product $A\widetilde{\otimes}B$, and by
    Proposition~5.4 the canonical tensor product
    $A\odot B$ is a composite in the sense of BGW Definition~4.1.
    Applying these to $A=h_3(\mathbb{C})$ or $h_3(\mathbb{H})$ with
    any non-trivial special EJA $B$ exhibits both as non-trivial
    Jordan-tensor factors.
  \item \textbf{$\hthree$ satisfies all three.}
    \begin{enumerate}[label=(\alph*)]
      \item (i) $\mathbb{K}=\mathbb{O}$: rank-1 idempotents in
        $\hthree$ have Peirce 0-space $V_0(E)\cong h_2(\mathbb{O})$
        (Proposition~\ref{prop:peirce-spaces}).
      \item (ii) $\mathbb{O}$ admits complex structures $u\in S^6
        \subset\mathrm{Im}\,\mathbb{O}$, so
        $h_2(\mathbb{C}_u)\hookrightarrow h_2(\mathbb{O})$
        (Lemma~\ref{lem:bottleneck}, proved below).
      \item (iii) $\hthree$ is the unique non-special simple FRJA
        by Zelmanov's classification~\cite{Zelmanov1979}; in
        particular it is not a JC-algebra.  By BGW
        Proposition~4.14~\cite{BGW2020}: \emph{if $A$ contains an
        exceptional ideal and $B$ contains a non-trivial ideal,
        there exists no composite $AB$ in the sense of BGW
        Definition~4.1.}  Equivalently~\cite[remark following
        Prop.~4.14]{BGW2020}, if $A$ is exceptional and any
        composite $AB$ exists, then $B$ must be a direct sum of
        one-dimensional Jordan algebras (a classical system).
        Hence $\hthree$ admits no non-trivial Jordan-tensor
        composite and therefore is not a non-trivial Jordan-tensor
        factor in the BGW category.
    \end{enumerate}
\end{enumerate}
By 1--5, $\hthree$ is the unique observer-grade JA basin up to
isomorphism.
\end{proof}

\begin{remark}[Scope of Theorem~\ref{thm:basin}; relation to Gap~A]
\label{rem:basin-scope}
Three categorical points are worth flagging.

\textbf{(a) ``Composite'' is BGW-Jordan-monoidal.}  Condition~(iii)
references the BGW Jordan-monoidal category~\cite{BGW2020}, which is
the natural categorical setting for formally real Jordan composites.
Other monoidal structures on Jordan-like categories may give other
notions of ``composite'' and other lists of basins;
Theorem~\ref{thm:basin} is sharp within BGW-Jordan-composability.

\textbf{(b) ``Universe'' is anthropic.}  The restriction to
\emph{observer-grade} basins is motivated by ``we are observers, so
we live in an observer-grade basin.''  Other simple finite-dimensional
formally real Jordan algebras (the spin factors and $h_n(\mathbb{K})$
for $\mathbb{K}\in\{\mathbb{R},\mathbb{C},\mathbb{H}\}$, $n\geq 3$)
exist as consistent mathematical structures; the program does not
claim they fail to host \emph{any} structure, only that they fail to
be observer-grade JA basins in the sense of
Definition~\ref{def:basin}.  Under Tegmark's Level~IV mathematical
realism~\cite{Tegmark2008} those alternatives exist and host whatever
structure their algebra permits; the program selects $\hthree$ as
\emph{ours} by anthropic restriction, not by uniqueness across all
possible worlds.  This is what was previously labeled Gap~A; the
present formulation makes the L4 step explicit and external to the
proof, rather than embedded in a prose argument.

\textbf{(c) Peirce sublocalization is not BGW-composition.}
Condition~(iii) refers to BGW Jordan-tensor factorization; the
arrow $A\to V_0(E)$ in this paper is a Peirce sublocalization
(Section~\ref{sec:peirce}).  These are different relations.  In
particular, $V_0(E) = h_2(\mathbb{O})$ is a Peirce 0-subspace of
$\hthree$ but not a BGW-Jordan factor of it; no contradiction with
$\hthree$ satisfying~(iii) arises.
\end{remark}

%----------------------------------------------------------------------
\subsection{Peirce decomposition under \texorpdfstring{$E_{11}$}{E11}}\label{sec:peirce}
%----------------------------------------------------------------------

The observer selects a rank-1 idempotent
\begin{equation}\label{eq:E11}
  e = E_{11} =
  \begin{pmatrix}
    1 & 0 & 0 \\
    0 & 0 & 0 \\
    0 & 0 & 0
  \end{pmatrix}.
\end{equation}

\begin{remark}[Gap~B, step~1]\label{rem:gap-B1}
The self-modeling framework requires that the universe contains a
faithful self-model---a subsystem whose states separate the states of
the whole.  Paper~5 proves that any such subsystem has
$M_n(\mathbb{C})^{\mathrm{sa}}$ state spaces, i.e.\ it is a
C*-algebra observer.  An observer of this type, embedded in $\hthree$,
occupies an idempotent whose Peirce 1-space $\Vone$ is its address
within the algebra (see Remark~\ref{rem:observer-universe}).  The
idempotent must be rank-1: a rank-2 idempotent would give
$\Vone \cong h_2(\mathbb{O})$ (10-dimensional), but the observer
is a C*-\emph{subalgebra} whose address is a single slot, not a
composite Jordan system.  Only rank-1 gives
$\Vone \cong \mathbb{R}$, the minimal address for a single
observer.
All rank-1 idempotents in $\hthree$ are conjugate under
$\Ffour$ (the orbit is the octonionic projective plane
$\mathbb{O}P^2$, $\dim = 16$).  The observer's existence
breaks $\Ffour \to \Spin{9}$; \emph{which} idempotent is selected
corresponds to the observer's ``viewpoint'' and is physically
irrelevant by $\Ffour$-conjugacy.  The breaking is forced by the
self-modeling definition (which requires an observer); the direction
is gauge.
\end{remark}

\begin{remark}[Observer-universe relation]\label{rem:observer-universe}
Paper~5 shows that a self-modeling system's state space is
$M_n(\mathbb{C})^{sa}$, which excludes $\hthree$ as an observer
algebra.  This is \emph{not} a contradiction: the observer is not
the universe.  The observer is \emph{internal} to $\hthree$ via the
Peirce decomposition: $\Vone \cong \mathbb{R}$ is the observer's
slot, $\Vhalf$ is the interaction space, and $\Vzero$ is the
complement.  The Peirce decomposition is the Jordan-algebraic analog
of a subsystem decomposition---not a tensor-product factorization
(which $\hthree$ does not admit, by non-composability) but a
spectral decomposition under the observer's idempotent.  The
observer's \emph{internal} algebra (its state space as a
self-modeling system) is C* by Paper~5; the \emph{universe's}
algebra is the non-composable $\hthree$.  The observer accesses the
universe through the Peirce interface $\Vhalf$, not through tensor
factors.

The dimensional relationship deserves emphasis.  The Peirce 1-space
$\Vone = \mathbb{R}\cdot E_{11}$ is the observer's \emph{address}
within $\hthree$---the eigenvalue-$1$ slot of the idempotent it
occupies---not the observer's internal state space.  The observer's
internal C*-algebra must be rich enough to faithfully represent the
Peirce operators $T_a \in \mathrm{End}(\Vhalf)$, which generate
$M_{16}(\mathbb{R})$.  Paper~5 forces this internal algebra to be
$M_n(\mathbb{C})^{\mathrm{sa}}$ for $n \geq 16$; the observer's
\emph{model} of the Peirce interface lives in
$M_{16}(\mathbb{C}) \subset M_n(\mathbb{C})$.  The 1-dimensional
Peirce slot and the $16^2$-dimensional measurement algebra are not
in conflict: the former is where the observer \emph{sits} in
$\hthree$, the latter is how it \emph{represents} the interface
$\Vhalf$ internally.  This is the standard situation for a subsystem
probing its environment: the subsystem's address is low-dimensional,
but its internal model of the environment is as rich as the
information accessible through the interface requires.
\end{remark}

The Peirce decomposition under $e$ partitions $\hthree$ into the
eigenspaces of the multiplication operator $L_e(X) = e \circ X$:
\begin{equation}\label{eq:peirce-decomp}
  \hthree \;=\; \Vone \;\oplus\; \Vhalf \;\oplus\; \Vzero,
  \qquad
  V_\lambda \;=\; \bigl\{X \in \hthree : e \circ X = \lambda X\bigr\}.
\end{equation}

\noindent\textbf{Explicit computation of $L_e(X)$.}  Writing
$E_{11}\, X + X\, E_{11}$ entry by entry and dividing by~$2$:
\begin{equation}\label{eq:L_eX}
  e \circ X \;=\;
  \frac{1}{2}
  \begin{pmatrix}
    2\alpha & \bar{x}_3 & x_2 \\[2pt]
    x_3 & 0 & 0 \\[2pt]
    \bar{x}_2 & 0 & 0
  \end{pmatrix}.
\end{equation}

Reading off the eigenvalue conditions:

\begin{proposition}[Peirce spaces]\label{prop:peirce-spaces}
The Peirce decomposition of $\hthree$ under $E_{11}$ is:
\begin{align}
  \Vone
    &= \bigl\{\alpha E_{11} : \alpha \in \mathbb{R}\bigr\}
    \cong \mathbb{R},
    & \dim(\Vone) &= 1,
    \label{eq:V1} \\
  \Vhalf
    &= \left\{\!
    \begin{pmatrix}
      0 & \bar{x}_3 & x_2 \\[2pt]
      x_3 & 0 & 0 \\[2pt]
      \bar{x}_2 & 0 & 0
    \end{pmatrix}
    : x_2,x_3 \in \mathbb{O}\right\}
    \cong \mathbb{O}^2,
    & \dim(\Vhalf) &= 16,
    \label{eq:V12} \\
  \Vzero
    &= \left\{\!
    \begin{pmatrix}
      0 & 0 & 0 \\[2pt]
      0 & \beta & \bar{x}_1 \\[2pt]
      0 & x_1 & \gamma
    \end{pmatrix}
    : \beta,\gamma\in\mathbb{R},\; x_1\in\mathbb{O}\right\}
    \cong h_2(\mathbb{O}),
    & \dim(\Vzero) &= 10.
    \label{eq:V0}
\end{align}
\end{proposition}

\noindent\textbf{Dimension check:}
$1 + 16 + 10 = 27 = \dim(\hthree)$.\ \checkmark

\medskip
\noindent\textbf{Stabilizer and Spin(9) representations.}
The stabilizer of $E_{11}$ in $\Ffour = \mathrm{Aut}(\hthree)$
($\dim\,\Ffour = 52$)
is~\cite{Baez2002,Yokota2009}
\begin{equation}\label{eq:stab-spin9}
  \mathrm{Stab}_{\Ffour}(E_{11}) \;=\; \Spin{9},
  \qquad
  \dim\,\Spin{9} = 36 = \tbinom{9}{2}.
\end{equation}
Coset check: $52 - 36 = 16 = \dim(\Vhalf)$,
consistent with $\Ffour/\Spin{9} \cong \mathbb{O}P^2$.

Under $\Spin{9}$ the Peirce spaces carry the representations:
\begin{center}
\renewcommand{\arraystretch}{1.2}
\begin{tabular}{@{}cccc@{}}
  \toprule
  Peirce space & Identification & $\dim_\mathbb{R}$ & $\Spin{9}$ rep \\
  \midrule
  $\Vone$ & $\mathbb{R}\cdot E_{11}$ & $1$ & trivial $\mathbf{1}$ \\
  $\Vhalf$ & $\mathbb{O}^2$ & $16$ & spinor $S_9 = \mathbf{16}$ \\
  $\Vzero$ & $h_2(\mathbb{O})$ & $10$ & $\mathbf{9}\oplus\mathbf{1}$ \\
  \midrule
  Total & $\hthree$ & $27$ & $\mathbf{1}\oplus\mathbf{16}\oplus\mathbf{9}\oplus\mathbf{1}$ \\
  \bottomrule
\end{tabular}
\end{center}

The identification $\Vhalf = \mathbb{O}^2 = S_9$
(the unique real 16-dimensional irreducible spinor of $\Spin{9}$)
follows from
$\dim_\mathbb{R}(S_9) = 2^{[9/2]} = 2^4 = 16$
and the fact that $\Spin{9}$ acts on the tangent space
$T_{E_{11}}(\mathbb{O}P^2) \cong \mathbb{O}^2$ via its
spinor representation~\cite{Baez2002,Yokota2009}.
The space $\Vzero = h_2(\mathbb{O})$ decomposes into the
traceless part (the 9-dimensional vector representation of
$\SO{9} = \Spin{9}/\mathbb{Z}_2$) plus the trace (a trivial
$\mathbf{1}$).

\medskip
\noindent\textbf{Peirce multiplication operators.}
For each unit vector $a$ in the 9-dimensional vector representation
(the traceless part of $\Vzero = h_2(\mathbb{O})$), the
\emph{Peirce multiplication operator} is
\begin{equation}\label{eq:T_a-def}
  T_a \colon \Vhalf \to \Vhalf,
  \qquad
  T_a(v) = (a \circ v)_{\Vhalf},
\end{equation}
where the subscript $\Vhalf$ denotes projection to the
Peirce half-space.  These are self-adjoint with respect to the
trace inner product on $\hthree$.  For an orthonormal basis
$\{a_1,\ldots,a_9\}$ of the vector representation,
\begin{equation}\label{eq:T_a-clifford}
  \{T_{a_i}, T_{a_j}\} = \tfrac{1}{2}\delta_{ij}\,I_{16},
\end{equation}
so $T_a^2 = \tfrac{1}{4}I_{16}$ and the rescaled operators
$\gamma_i = 2T_{a_i}$ satisfy the standard Clifford relation
$\{\gamma_i,\gamma_j\} = 2\delta_{ij}$.  The image of the
Clifford algebra $\mathrm{Cl}(9,0)$ in
$\mathrm{End}(\Vhalf) = M_{16}(\mathbb{R})$ is
$M_{16}(\mathbb{R})$ itself: the abstract algebra
$\mathrm{Cl}(9,0) \cong M_{16}(\mathbb{R})
\oplus M_{16}(\mathbb{R})$ has two simple summands, and the spinor
representation $S_9$ is the irreducible module for one summand.

%----------------------------------------------------------------------
\subsection{C*-bottleneck universality}\label{sec:bottleneck}
%----------------------------------------------------------------------

A C*-observer accessing $\hthree$ probes the algebra through positive
unital projections (conditional expectations) onto C*-subalgebras.
A natural reviewer question is whether the complexification of
Section~\ref{sec:complexification-arg} below depends on a free choice
of C*-target: ``why $h_3(\mathbb{C}_u)$ rather than another convenient
slice?''  Lemma~\ref{lem:bottleneck} answers that the family of
admissible C*-targets is structurally determined: up to
$\Ffour$-action, the maximal-range targets are exactly the algebras
$h_3(\mathbb{C}_u)$ indexed by $u\in S^6\subset\mathrm{Im}\,\mathbb{O}$.
The selection of $u$ inside $S^6$ is then necessary but
direction-irrelevant (Gap~B2; cf.~Section~\ref{sec:gap-register}).

\begin{lemma}[C*-bottleneck universality]\label{lem:bottleneck}
Let $\mathcal{E}(\hthree)$ denote the family of positive unital
idempotents (conditional expectations)
$\mathcal{E}\colon\hthree\to A$ with $A\subseteq\hthree$ a Jordan
subalgebra (the embedding $A\hookrightarrow\hthree$ preserves the
Jordan product) that is isomorphic to the self-adjoint part of a
finite-dimensional complex C*-algebra.  Then:
\begin{enumerate}[label=(\roman*)]
  \item For every $\mathcal{E}\in\mathcal{E}(\hthree)$, the range
    $A=\mathcal{E}(\hthree)$ is a real Jordan subalgebra of
    $\hthree$ of the form
    $A\cong\bigoplus_{k=1}^{r} M_{n_k}(\mathbb{C})^{sa}$ with
    $\sum_k n_k\leq 3$ and (consequently) each $n_k\leq 3$.
  \item The maximal-range elements of $\mathcal{E}(\hthree)$
    (those with $\dim_{\mathbb{R}} A$ maximal) have
    $A\cong M_3(\mathbb{C})^{sa}\cong h_3(\mathbb{C}_u)$ for some
    $u\in S^6\subset\mathrm{Im}\,\mathbb{O}$, where
    $\mathbb{C}_u=\mathbb{R}\oplus\mathbb{R}u$ is the complex line
    through~$u$.  The family $\{h_3(\mathbb{C}_u):u\in S^6\}$ forms
    a single $\Ffour$-orbit.
  \item For any rank-1 idempotent $E\in h_3(\mathbb{C}_u)$, the
    Peirce 0-space within the maximal C*-target satisfies
    $V_0(E)\cap h_3(\mathbb{C}_u)\cong h_2(\mathbb{C}_u)\cong
    M_2(\mathbb{C})^{sa}$.
\end{enumerate}
\end{lemma}

\begin{proof}
\textbf{(i)}~Effros-St{\o}rmer~\cite{EffrosStormer1979}: the range
of a positive unital idempotent on a JB-algebra is a JB-subalgebra.
In finite dimensions over $\mathbb{R}$ this is just a Jordan
subalgebra; by the hypothesis on $\mathcal{E}$ the embedding is
moreover Jordan-product-preserving, so $A\subseteq\hthree$ as a
Jordan subalgebra.  The hypothesis that $A$ is isomorphic to the
self-adjoint part of a finite-dimensional complex C*-algebra,
combined with the Artin--Wedderburn classification, gives
$A\cong\bigoplus_{k=1}^{r} M_{n_k}(\mathbb{C})^{sa}$ for some
positive integers $n_1,\ldots,n_r$.  The Jordan rank of
$M_{n_k}(\mathbb{C})^{sa}$ is $n_k$ (the maximum number of
mutually orthogonal idempotents), so the rank of $A$ is
$\sum_k n_k$.  Since $A$ is a Jordan subalgebra of $\hthree$ and
Jordan rank decreases under Jordan subalgebra inclusion, we have
$\sum_k n_k\leq\mathrm{rank}(\hthree)=3$.  In particular each
$n_k\leq 3$.

\textbf{(ii)}~Maximizing $\dim_{\mathbb{R}} A$ subject to
$\sum_k n_k\leq 3$ and each summand $M_{n_k}(\mathbb{C})^{sa}$
having $\dim_{\mathbb{R}}=n_k^2$, the unique maximizer is
$r=1$, $n_1=3$, giving $\dim_{\mathbb{R}} A=9$.  (The other
options---$M_2\oplus M_1$, $M_1\oplus M_1\oplus M_1$, $M_2$, $M_1$,
etc.---all give $\dim_{\mathbb{R}} A\leq 5$.)  So a maximal
$A\cong M_3(\mathbb{C})^{sa}$.  We now identify the embedding into
$\hthree$.  Fixing any $u\in\mathbb{O}$ with $u^2=-1$ defines the
associative subfield $\mathbb{C}_u=\mathbb{R}\oplus\mathbb{R}u
\subset\mathbb{O}$, and
\[
  h_3(\mathbb{C}_u) \;:=\; \{X\in\hthree : X_{ij}\in\mathbb{C}_u
  \text{ for all } i,j\} \;\cong\; M_3(\mathbb{C})^{sa}
\]
is a Jordan subalgebra of $\hthree$ of dimension~$9$.  By the
Baez classification of Jordan subalgebras of
$\hthree$~\cite[\S 3.4 \& \S 4.4]{Baez2002}, the rank-3
Jordan subalgebras of $\hthree$ isomorphic to $M_n(\mathbb{C})^{sa}$
are exactly $h_3(\mathbb{C}_u)$ for $u\in S^6
\subset\mathrm{Im}\,\mathbb{O}$ (the rank-3 quaternionic and real
forms $h_3(\mathbb{H}),h_3(\mathbb{R})\subset\hthree$ are not
isomorphic to any $M_n(\mathbb{C})^{sa}$ and so do not occur in
$\mathcal{E}(\hthree)$).  The automorphism group
$G_2=\Aut(\mathbb{O})\subset\Ffour$ acts transitively on the
imaginary unit sphere ($S^6\cong G_2/\SU{3}$), so
$\{h_3(\mathbb{C}_u)\}_{u\in S^6}$ forms a single $\Ffour$-orbit.

\textbf{(iii)}~Within $h_3(\mathbb{C}_u)\cong M_3(\mathbb{C})^{sa}$,
a rank-1 idempotent $E$ has the standard Peirce decomposition
$h_3(\mathbb{C}_u) = V_1(E)\oplus V_{1/2}(E)\oplus V_0(E)$
with $V_1(E)\cong\mathbb{R}$, $V_{1/2}(E)\cong\mathbb{C}_u^2$,
and $V_0(E)\cong h_2(\mathbb{C}_u)\cong M_2(\mathbb{C})^{sa}$.
Hence $V_0(E)\cap h_3(\mathbb{C}_u)\cong h_2(\mathbb{C}_u)$.
\end{proof}

\begin{remark}[Lorentzian signature on
$h_2(\mathbb{C}_u)$]\label{rem:minkowski}
The space $h_2(\mathbb{C}_u)\cong M_2(\mathbb{C})^{sa}$ carries the
real quadratic form $X\mapsto -\det(X)$, which is of signature
$(1,3)$ when $X$ is written in the basis
$(\sigma_0,\sigma_1,\sigma_2,\sigma_3)$ of $2\times 2$ Hermitian
matrices.  This is the standard Minkowski metric on Hermitian
$2\times 2$ matrices.  Lemma~\ref{lem:bottleneck}~(iii) records the
spectator complement of the C*-bottleneck as $h_2(\mathbb{C}_u)$ and
this remark notes that it carries Lorentzian signature canonically;
the use of the bottleneck for spacetime identification is downstream
work and does not affect the chirality derivation in this paper.
\end{remark}

\begin{remark}[Selection vs.\ forcing]\label{rem:sel-vs-force}
Lemma~\ref{lem:bottleneck} establishes that the menu of maximal
C*-targets is the $\Ffour$-orbit
$\{h_3(\mathbb{C}_u):u\in S^6\}$, no more and no less.  The selection
of a particular $u\in S^6$ (Gap~B2; see
Section~\ref{sec:gap-register}) is required: Schur's lemma
($\mathrm{End}_{\Spin{9}}(S_9)=\mathbb{R}$ since $9\bmod 8=1$)
forbids equivariant complexification, so no canonical $u$ exists.
The selection produces $\Ffour$-conjugate physics, so the direction
is gauge.  The complexification of $\Vhalf$ established by
Theorem~\ref{thm:complexification} below is the dynamical realization
of one such selection: when the observer's continuous functional
calculus is applied to a specific Peirce operator $T_a$, the resulting
imaginary structure pins down an associated $u\in S^6$, and the
universal-property statement here certifies that this is the maximal
admissible target rather than a convenient one.
\end{remark}

%----------------------------------------------------------------------
\subsection{Complexification from C*-observer nature}\label{sec:complexification-arg}
%----------------------------------------------------------------------

The central claim of Part~A is that \emph{complexification is a
consequence of the observer's algebraic nature, not assumed}.  We
prove this via the continuous functional calculus of the observer's
C*-algebra, applied to the Peirce operators.

\medskip
\noindent\textbf{Step~1: Self-modeling forces a C*-algebra observer.}
Paper~5~\cite{Paper5} defines a \emph{self-modeling system} as a
finite-dimensional order-unit space $(V, e)$ containing a
sub-order-unit space $(\hat{V}, \hat{e})$ satisfying four axioms:
(A1)~the states of $\hat{V}$ separate the states of $V$
(faithfulness); (A2)~$\hat{V}$ is updated through the boundary
between the subsystem and its complement; (A3)~the sequential product
$a\,\&\,b = \sqrt{a}\,b\,\sqrt{a}$ defines the measurement
composition; (A4)~$\hat{V}$ is closed under this product.
Paper~5 proves that these axioms force $V \cong
M_n(\mathbb{C})^{\mathrm{sa}}$ for some $n$---the self-adjoint
part of a matrix C*-algebra.

The square root in the sequential product is computed via the
\emph{continuous functional calculus} (CFC) of the ambient C*-algebra
$M_n(\mathbb{C})$.  For a self-adjoint element $x$ with spectral
decomposition $x = \sum_i \lambda_i P_i$, this calculus defines
$f(x) = \sum_i f(\lambda_i)\,P_i$ for any continuous
$f\colon\sigma(x)\to\mathbb{C}$.  This is the \emph{scalar}
functional calculus: it applies $f$ to each eigenvalue
independently and is uniquely determined by the C*-algebra
structure.  The distinction between this scalar calculus and general
matrix-function constructions is critical for
Section~\ref{sec:complexification-arg}
(see Remark~\ref{rem:scalar-vs-matrix-sqrt}).

\medskip
\noindent\textbf{Step~2: Peirce operators have negative eigenvalues.}
The Peirce multiplication operators $T_a$ (Section~\ref{sec:peirce})
satisfy $T_a^2 = \tfrac{1}{4}I_{16}$, so each $T_a$ has spectrum
$\{+\tfrac{1}{2}, -\tfrac{1}{2}\}$.  The eigenvalue $-\tfrac{1}{2}$
is \emph{negative}.  In a real algebra, $\sqrt{-\tfrac{1}{2}}$ is
undefined.  But the observer is $M_n(\mathbb{C})^{\mathrm{sa}}$
(Step~1), so the continuous functional calculus of $M_n(\mathbb{C})$
applies the principal branch:
\begin{equation}\label{eq:sqrt-neg}
  \sqrt{-\tfrac{1}{2}} = \frac{i}{\sqrt{2}}.
\end{equation}
This is the principal branch value
($\arg(\sqrt{z})\in(-\pi/2,\pi/2]$); the sign of~$i$ depends
on branch convention (see below), but the exit from $\mathbb{R}$
does not.

\medskip
\noindent\textbf{Step~3: The CFC computation exits $M_{16}(\mathbb{R})$.}
The observer faithfully represents the Peirce operators as self-adjoint
elements of $M_n(\mathbb{C})$ via the standard inclusion
$\iota: M_{16}(\mathbb{R}) \hookrightarrow M_{16}(\mathbb{C})$
($n \geq 16$).  The spectral decomposition of $T_a$ gives
\begin{equation}\label{eq:sqrt-Ta}
  \sqrt{T_a} = \frac{1}{\sqrt{2}}(P_+ + i\,P_-),
\end{equation}
where $P_\pm$ are the spectral projections onto the $\pm\tfrac{1}{2}$
eigenspaces.

\begin{lemma}[CFC well-definedness]\label{lem:sp-domain}
Let $T_a, T_b$ be self-adjoint Peirce operators, faithfully represented
in the observer's C*-algebra $M_n(\mathbb{C})$ via
$\iota\colon M_{16}(\mathbb{R})\hookrightarrow M_n(\mathbb{C})$.
Then:
\begin{enumerate}
  \item[\textup{(a)}] $\sqrt{\iota(T_a)}\,\iota(T_b)\,\sqrt{\iota(T_a)}$
    is well-defined via the continuous functional calculus of
    $M_n(\mathbb{C})$.
  \item[\textup{(b)}] For effects $a \in [0,I]$, this operation coincides
    with the sequential product $a\,\&\,b = \sqrt{a}\,b\,\sqrt{a}$
    derived in Paper~5~\cite{Paper5}.
\end{enumerate}
\end{lemma}

\begin{proof}
(a)~For any self-adjoint $x \in M_n(\mathbb{C})^{sa}$ with spectral
decomposition $x = \sum_i \lambda_i P_i$, the continuous functional
calculus defines $f(x) = \sum_i f(\lambda_i)\,P_i$ for every
continuous $f\colon\sigma(x)\to\mathbb{C}$.  Applying this to
$f(z) = z^{1/2}$ on $\sigma(\iota(T_a)) = \{+\tfrac{1}{2},
-\tfrac{1}{2}\}$ yields $\sqrt{\iota(T_a)} \in M_n(\mathbb{C})$
(the branch choice on $-\tfrac{1}{2}$ is addressed below
Theorem~\ref{thm:complexification}).  The triple product
$\sqrt{\iota(T_a)}\,\iota(T_b)\,\sqrt{\iota(T_a)}$ is then
a well-defined element of $M_n(\mathbb{C})$.
(b)~Paper~5 derives the L\"uders form
$a\,\&\,b = \sqrt{a}\,b\,\sqrt{a}$ for the sequential product
on effects, where $\sqrt{a}$ is the unique positive square root---itself
a CFC element.  For effects, the positive square root and the CFC
square root coincide, so the two computations are identical.
\end{proof}

\begin{theorem}[Observer-induced complexification]\label{thm:complexification}
Let $T_a, T_b$ be anticommuting Peirce operators
($\{T_a, T_b\} = 0$, $a \neq b$) in
$\mathrm{End}(\Vhalf) = M_{16}(\mathbb{R})$, and let
$\iota\colon M_{16}(\mathbb{R})\hookrightarrow M_{16}(\mathbb{C})
\subset M_n(\mathbb{C})$ be the faithful representation of the Peirce
operators inside the observer's C*-algebra (Step~3).  Then the
operation $\sqrt{T_a}\,T_b\,\sqrt{T_a}$
(Lemma~\ref{lem:sp-domain}), computed via the continuous functional
calculus of $M_n(\mathbb{C})$, gives
\begin{equation}\label{eq:sp-complex}
  \sqrt{\iota(T_a)}\;\iota(T_b)\;\sqrt{\iota(T_a)}
  = \frac{i}{2}\,\iota(T_b).
\end{equation}
We write $\sqrt{T_a}$ for $\sqrt{\iota(T_a)}$ when the representation
context is clear.
\end{theorem}

\begin{proof}
Expand using~\eqref{eq:sqrt-Ta}:
\begin{equation}\label{eq:sp-expand}
  \sqrt{T_a}\; T_b\; \sqrt{T_a}
  = \tfrac{1}{2}(P_+ + iP_-)\,T_b\,(P_+ + iP_-)
  = \tfrac{1}{2}\bigl(P_+ T_b P_+ + i\,P_+ T_b P_-
    + i\,P_- T_b P_+ - P_- T_b P_-\bigr).
\end{equation}
From $\{T_a, T_b\} = 0$ and $P_+ T_a = \tfrac{1}{2}P_+$, it follows
that $P_+ T_b P_+ = 0$ and $P_- T_b P_- = 0$ (anticommuting Hermitian
matrices have purely off-diagonal structure in each other's eigenbasis).
Therefore $T_b = P_+ T_b P_- + P_- T_b P_+$, and
\begin{equation}
  \sqrt{T_a}\; T_b\; \sqrt{T_a}
  = \tfrac{i}{2}(P_+ T_b P_- + P_- T_b P_+) = \tfrac{i}{2}\,T_b.
  \qedhere
\end{equation}
\end{proof}

\noindent
The specific sign of~$i$ in~\eqref{eq:sp-complex} depends on the
branch convention (replacing $i$ by $-i$ corresponds to choosing the
other continuous branch of $\sqrt{z}$ on
$\{+\tfrac{1}{2},-\tfrac{1}{2}\}$), but the exit from $\mathbb{R}$
to $\mathbb{C}$ does not: every square-root choice for eigenvalue
$-\tfrac{1}{2}$ produces a non-real value, so the measurement
algebra exits $M_{16}(\mathbb{R})$ regardless of convention.

The factor of~$i$ is determined by the functional calculus, not
conventional.  Paper~5 establishes
that a self-modeling observer's measurement composition is the
sequential product $a\,\&\,b = \sqrt{a}\,b\,\sqrt{a}$ for effects,
computed via the continuous functional calculus of its C*-algebra
$M_n(\mathbb{C})$.  The same functional calculus applies to all
self-adjoint elements of $M_n(\mathbb{C})$, including the Peirce
operators $T_a$ (Lemma~\ref{lem:sp-domain}).  When the observer
probes the Peirce half-space
via the operators $T_a$---which are the natural observables of
the Peirce interface (Eq.~\eqref{eq:T_a-def})---the CFC computation
$\sqrt{T_a}\,T_b\,\sqrt{T_a}$ necessarily involves $\sqrt{T_a}$.  Since $T_a$ has
eigenvalue $-\tfrac{1}{2} < 0$, the real functional calculus
cannot produce $\sqrt{T_a}$; only the complex functional calculus
applies (Eq.~\eqref{eq:sqrt-neg}).  By contrast, the
effect operators $E_a = \tfrac{1}{2}(I_{16} + 2T_a)$ have spectrum
$\{0, 1\} \subset [0, 1]$, so $\sqrt{E_a}$ is real-valued and the
sequential product $E_a\,\&\,E_b = \tfrac{1}{2}E_a$ remains entirely
real.  The complexification arises specifically from the negative
eigenvalues of the Peirce operators, not from a choice to extend
beyond the effect domain.

\begin{remark}[Scalar functional calculus vs.\ matrix square roots]
\label{rem:scalar-vs-matrix-sqrt}
A potential objection is that $T_a$ admits a \emph{real} matrix
square root: since the eigenvalue $-\tfrac{1}{2}$ has even
multiplicity ($8$), one can pair eigenspaces into $2\times 2$ blocks
and use real rotation matrices $J$ ($J^2 = -I_2$) to construct
$S\in M_{16}(\mathbb{R})$ with $S^2 = T_a$
(cf.~\cite{Higham2008}, Thm.~1.29).

This objection conflates two distinct notions of ``square root.''
The C*-algebra continuous functional calculus (Step~1) is the
\emph{scalar} functional calculus: it defines
$\sqrt{T_a} = \sum_i \sqrt{\lambda_i}\,P_i$, applying the
square-root function to each eigenvalue \emph{independently}.
For $\lambda = -\tfrac{1}{2}$, the equation $z^2 = -\tfrac{1}{2}$
has no real solution, so the functional calculus square root is
necessarily complex.  The real matrix square root $S$ described
above is \emph{not} a functional calculus element: it cannot be
written as $g(T_a) = \sum_i g(\lambda_i)\,P_i$ for any scalar
function $g$, because it mixes eigenspaces via $2\times 2$
rotation blocks rather than acting on each eigenspace by a scalar.

Paper~5 establishes that the self-modeling sequential product is
computed via the C*-algebra continuous functional calculus.  This is
not a choice among available square roots but a consequence of the
C*-algebraic structure: the functional calculus is uniquely
determined, and it is the scalar one.  The Higham-type real matrix
square roots, while mathematically valid, are not the square roots
that the self-modeling framework selects.
\end{remark}

\begin{remark}[Why the observer measures $T_a$, not just effects]
\label{rem:sp-domain}
A natural objection is that the Gudder--Greechie sequential
product~\cite{GudderGreechie2002} is defined for effects (spectrum
in $[0,1]$), while $T_a$ has eigenvalue $-\tfrac{1}{2}$.
Lemma~\ref{lem:sp-domain} establishes that
$\sqrt{T_a}\,T_b\,\sqrt{T_a}$ is well-defined in the observer's
C*-algebra via the continuous functional calculus, and that this
operation coincides with the sequential product on effects.
Here we address why this computation is physically relevant for
the Peirce operators, which lie outside the effect interval.

\emph{$T_a$ is the algebraic primitive; $E_a$ is a derived
rescaling.}\enspace
The Peirce operator $T_a(v) = (a\circ v)_{\Vhalf}$ \emph{is} the
Jordan product restricted to $\Vhalf$---the fundamental algebraic
operation at the observer's interface with the universe.
The effect $E_a = \tfrac{1}{2}(I + 2T_a)$ is an affine
transformation of $T_a$ into the unit interval---a derived object,
not the fundamental one.  The objection that ``the observer should
use $E_a$ instead of $T_a$'' reverses the logical order.

\emph{$T_a$ and $E_a$ are different observables.}\enspace
$T_a$ has eigenvalues $\pm\tfrac{1}{2}$; $E_a$ has eigenvalues
$\{0,1\}$.  They represent different measurement outcomes, and the
CFC computations $\sqrt{T_a}\,T_b\,\sqrt{T_a}$ and
$\sqrt{E_a}\,E_b\,\sqrt{E_a}$ encode different physical
compositions.  Shifting the spectrum is not a gauge freedom---it
changes what is being measured.

\emph{$T_a$ is in the algebra necessarily.}\enspace
The Peirce operators generate
$\mathrm{End}(\Vhalf) = M_{16}(\mathbb{R})$: they are the
basis of the algebra of observable operations on the Peirce
interface.  Any faithful representation of this algebra in the
observer's C*-algebra $M_n(\mathbb{C})$ must include the $T_a$
as self-adjoint elements.  Once
$T_a \in M_n(\mathbb{C})^{\mathrm{sa}}$, the continuous
functional calculus defines $\sqrt{T_a}$ uniquely, and the
CFC computation $\sqrt{T_a}\,T_b\,\sqrt{T_a}$ is a
well-defined element of $M_n(\mathbb{C})$.  The complexification
is therefore a structural consequence of the algebra's
C*-nature---it follows from the \emph{existence} of $T_a$ in
the algebra, not from any particular measurement the observer
performs.
\end{remark}

\begin{remark}[Sequential product output and the generated algebra]
\label{rem:sp-output}
Theorem~\ref{thm:complexification} produces
$\sqrt{T_a}\,T_b\,\sqrt{T_a} = (i/2)\,T_b$, which is
anti-self-adjoint---not an observable in the standard sense.
This is not a defect of the argument but its operational content.

The CFC operation $\sqrt{T_a}\,(\cdot)\,\sqrt{T_a}$ acts within the observer's
C*-algebra $M_n(\mathbb{C})$; its closure generates a
\emph{subalgebra}, not merely a collection of self-adjoint elements.
The subalgebra generated by $\{T_a\}$ together with all iterated
CFC products is $M_{16}(\mathbb{C})$, not
$M_{16}(\mathbb{R})$.  The physical consequence is that the
observer's \emph{observable space}---the self-adjoint part of the
generated algebra---is $M_{16}(\mathbb{C})^{\mathrm{sa}}$, which is
strictly larger than $M_{16}(\mathbb{R})$.  Elements such as
$T_a + (i/2)\,T_b$ are self-adjoint in $M_{16}(\mathbb{C})$ but
have no counterpart in $M_{16}(\mathbb{R})$: they are new
observables that exist only because the algebra extends to
$M_{16}(\mathbb{C})$.

The complexification is therefore a statement about the
\emph{algebra of all composable measurements}, not about any single
measurement outcome.  The symmetries of the Peirce interface are
those of this generated algebra---$\Spin{10}$ acting on
$\mathbb{C}^{16}$---not those of the original real
representation.
\end{remark}

The algebra generated by the observer's CFC computations
$\sqrt{T_a}\,T_b\,\sqrt{T_a}$ on the
Peirce operators is therefore $M_{16}(\mathbb{C})$, not
$M_{16}(\mathbb{R})$: the $\mathbb{C}$-linear span of all monomials
in the $T_a$ and their CFC products has complex dimension~256,
equal to $\dim_\mathbb{C} M_{16}(\mathbb{C})$.  This is the image
of the complexified Clifford algebra $\mathrm{Cl}(9,\mathbb{C})
\cong M_{16}(\mathbb{C}) \oplus M_{16}(\mathbb{C})$ acting on
$\Vhalf^{\,\mathbb{C}} = \mathbb{C}^{16}$ via one of its two
simple summands (cf.\ the real case in
Eq.~\eqref{eq:T_a-clifford}).  The observer's effective description
of $\Vhalf$ is therefore
\begin{equation}\label{eq:complexification}
  \Vhalf^{\,\mathbb{C}}
  \;:=\;
  \Vhalf \otimes_\mathbb{R} \mathbb{C}\,,
  \qquad
  \dim_\mathbb{C}\!\bigl(\Vhalf^{\,\mathbb{C}}\bigr)
  = 16.
\end{equation}

\begin{remark}[Compatibility with impossibility theorems]
\label{rem:impossibility-compat}
No $\Spin{9}$-equivariant complex structure $J: \Vhalf \to \Vhalf$
with $J^2 = -\mathrm{Id}$ exists (Schur's lemma:
$\mathrm{End}_{\Spin{9}}(S_9) = \mathbb{R}$, since
$9 \bmod 8 = 1$ gives real type by Bott periodicity).  The
observer-induced complexification is \emph{non-equivariant}: it
depends on which $T_a$ the observer measures first.  Different
choices of $T_a$ produce different complex structures, corresponding
to different choices of $u \in S^6 \subset \mathrm{Im}(\mathbb{O})$
(Gap~B, step~2).  The correspondence is explicit: each Peirce
operator $T_a$ is indexed by a unit vector $a$ in the
9-dimensional traceless part of $\Vzero = h_2(\mathbb{O})$; for
$a = e_k \in \mathrm{Im}(\mathbb{O})$ ($k=1,\ldots,7$), the
complex structure induced on $\Vhalf = \mathbb{O}^2$ by the
CFC computation is left multiplication by $e_k$, giving $u = e_k$.
All seven such complex structures are $\Gfour$-conjugate (orbit
$= S^6$, dim~6).
The impossibility of equivariant complexification
is complemented, not circumvented, by the observer's external input.
\end{remark}

\begin{remark}[Complexification vs.\ complex structure]
\label{rem:C-vs-B2}
Theorem~\ref{thm:complexification} establishes that complexification
\emph{occurs} for any pair of anticommuting Peirce operators---no
choice of $u$ is required.  The choice of $u \in S^6$ (Gap~B,
step~2) determines \emph{which} complex structure the observer
obtains, not \emph{whether} complexification occurs.  These are
logically distinct: the former (Gap~C, that complexification \emph{occurs}) is closed conditional on the observer's $T_a$-probing premise (Remark~\ref{rem:sp-domain}; equivariant forcing is impossible by Schur)---not unconditionally
by Theorem~\ref{thm:complexification}; the latter (Gap~B, step~2)
is a necessary symmetry reduction forced by the combination of
Theorem~\ref{thm:complexification} and Schur's lemma
(Remark~\ref{rem:impossibility-compat}).  Since all
$u \in S^6$ are $\Gfour$-conjugate, the reduction is required but
the direction is physically irrelevant---a gauge choice
(cf.~\cite{FerraraEtAl2021} for an explicit $\Gfour\to\SU{3}$
model on $S^6$).

Crucially, B2 does not await external dynamics.  In the self-modeling
framework, the observer's \emph{only} dynamics is the sequential
product: the ``before $\to$ model $\to$ after'' cycle of
self-modeling measurement (Paper~5).  Moreover, the observer
\emph{must} probe $\Vhalf$: Paper~5 requires a faithful self-model
updated through the system's boundary, and $\Vhalf$ \emph{is} the
Peirce boundary---the only sector through which the observer's slot
$\Vone$ communicates with the complement $\Vzero$ (the Peirce
multiplication rules give $\Vone\circ\Vzero = 0$, so
$\Vhalf$ is the sole channel).  An observer that ignores $\Vhalf$
has no model of its environment and fails the faithfulness condition.

Given that $\Vhalf$ must be probed, the observer cannot do so
without applying $\sqrt{T_a}\,(\cdot)\,\sqrt{T_a}$ for
some specific~$T_a$; \emph{every} such application involves a
specific $T_a$ and thereby a specific $u$.  There is no state of the
self-modeling cycle in which $\Vhalf$ is accessed but no $T_a$ has
been selected: probing $\Vhalf$ and selecting a $T_a$ are the same
algebraic act.  That the observer must access $\Vhalf$ via
\emph{some} $T_a$ is forced; \emph{which} $T_a$---and hence which
$u \in S^6$---is not determined by the algebra, which is what makes
B2 a genuine necessary symmetry reduction rather than a
derived result.  The observer's C*-nature---expressed through the CFC computation
$\sqrt{T_a}\,(\cdot)\,\sqrt{T_a}$---is therefore both the
mechanism that produces complexification
(Theorem~\ref{thm:complexification}) and the act that selects the
complex structure.  No temporal ordering is required---the point is
structural: any instance of the self-modeling cycle that accesses
$\Vhalf$ necessarily involves a specific $T_a$, hence a
specific~$u$.
\end{remark}

\begin{remark}[Contrast with Boyle~\cite{Boyle2020}]\label{rem:boyle-contrast}
Boyle starts from the complexified algebra
$\hthreeC := \hthree\otimes_\mathbb{R}\mathbb{C}$ and
studies its structure group $\Esix$.  Our approach differs in that
we \emph{derive} the complexification from the observer's C*-nature:
the self-modeling requirement (Paper~5) forces the observer to be a
C*-algebra system, and the continuous functional calculus applied to the Peirce operators
necessarily exits $M_{16}(\mathbb{R})$ into $M_{16}(\mathbb{C})$
(Theorem~\ref{thm:complexification}).  The end result is the same
identification of $\Vhalf^{\,\mathbb{C}}$ as the Weyl spinor
$\Stenplus$, but the logical route is different: Boyle assumes
complexification; we derive it from self-modeling.
\end{remark}

%----------------------------------------------------------------------
\subsection{Spin(9) to Spin(10) and \texorpdfstring{$\Ffour$}{F4} to \texorpdfstring{$\Esix$}{E6}}\label{sec:spin-upgrade}
%----------------------------------------------------------------------

The complexification of $\Vhalf$ triggers a representation-theoretic
upgrade from $\Spin{9}$ to $\Spin{10}$.

\begin{proposition}[Spinor upgrade]\label{prop:spinor-upgrade}
The complexified Peirce half-space
$\Vhalf^{\,\mathbb{C}} = \mathbb{O}^2\otimes_\mathbb{R}\mathbb{C}$
carries the Weyl spinor representation $\Stenplus$ of $\Spin{10}$:
\begin{equation}\label{eq:branching}
  \Stenplus\big|_{\Spin{9}}
  \;\cong\;
  S_9 \otimes_\mathbb{R} \mathbb{C}
  \;=\;
  S_9^{\,\mathbb{C}}.
\end{equation}
\end{proposition}

\begin{proof}
$\Spin{10}$ has two inequivalent complex Weyl representations
$\Stenplus$ and $S_{10}^-$, each of
$\dim_\mathbb{C} = 2^{[10/2]-1} = 2^4 = 16$.
The standard branching rule for the embedding
$\Spin{9}\hookrightarrow\Spin{10}$ (where $\Spin{9}$ stabilizes a
unit vector in $\mathbb{R}^{10}$) gives
$\Stenplus|_{\Spin{9}} \cong S_9^{\,\mathbb{C}}$:
the Weyl spinor of $\Spin{10}$, restricted to $\Spin{9}$, is
the complexification of the real spinor $S_9$.  Equivalently,
$S_9$ is a real form of $\Stenplus$.

Since $\Vhalf = S_9$ (Proposition~\ref{prop:peirce-spaces}) and
$\Vhalf^{\,\mathbb{C}} = S_9^{\,\mathbb{C}}$, we conclude
$\Vhalf^{\,\mathbb{C}} = \Stenplus$.
\end{proof}

\noindent\textbf{Dimension check:}
$\dim_\mathbb{C}(\Stenplus) = 16 = \dim_\mathbb{R}(S_9)
= \dim_\mathbb{R}(\Vhalf)$.\ \checkmark

\begin{remark}[Why $\Spin{10}$, not just complex $\Spin{9}$]
\label{rem:why-spin10}
A natural objection is that complexifying $S_9$ gives a complex
$\Spin{9}$-module, not automatically a $\Spin{10}$-module.  The
upgrade is structural, not merely representation-theoretic.

\emph{Representation-theoretic uniqueness.}\enspace
The branching $\Stenplus|_{\Spin{9}} \cong S_9^{\,\mathbb{C}}$
is multiplicity-free: $S_9^{\,\mathbb{C}}$ is irreducible under
$\Spin{9}$.  By Schur's lemma, any $\Spin{10}$-module structure on
$S_9^{\,\mathbb{C}}$ that restricts to the given $\Spin{9}$ action is
unique (up to the $\Stenplus$ vs.\ $S_{10}^-$ choice resolved below).

\emph{Structural forcing.}\enspace
The decisive point is not representation compatibility but
the structure group of $\hthreeC$.
The complexification $\hthree \to \hthreeC$ upgrades the
automorphism group from $\Ffour$ to the structure group $\Esix$,
and correspondingly upgrades the stabilizer of $E_{11}$ from
$\Spin{9}$ to $\Spin{10}\times\mathrm{U}(1)$
(Eq.~\eqref{eq:stab-E6}).  The additional $9$ generators of
$\mathfrak{spin}(10)$ beyond $\mathfrak{spin}(9)$ are
structure-group elements of $\hthreeC$ that do not preserve the real
$\hthree$: they exist in $M_{16}(\mathbb{C})$ but not in
$M_{16}(\mathbb{R})$, and they are \emph{constructed} by the
complexification that Theorem~\ref{thm:complexification} establishes.
Thus $\Spin{10}$ is the symmetry of the observer's complexified
description of the Peirce interface---the stabilizer of $E_{11}$
within the structure group of the algebra the observer actually
works in---not merely a group whose representation ring happens
to contain~$S_9^{\,\mathbb{C}}$.

\emph{Why the upgrade is physical, not merely formal.}\enspace
The observer cannot ``ignore'' the complexification and retain
$\Spin{9}$ as the relevant symmetry group.
Theorem~\ref{thm:complexification} establishes that the
observer's measurement algebra is $M_{16}(\mathbb{C})$, not
$M_{16}(\mathbb{R})$: the CFC computations it performs
generate complex elements that do not lie in
$M_{16}(\mathbb{R})$.  The symmetry group of the observer's
description of $\Vhalf$ is therefore the stabilizer of $E_{11}$
in the structure group of $\hthreeC$ (which is $\Esix$), not in
the automorphism group of~$\hthree$ (which is $\Ffour$).  This
gives $\Spin{10}\times\mathrm{U}(1)$
(Eq.~\eqref{eq:stab-E6}), not $\Spin{9}$
(Eq.~\eqref{eq:stab-spin9}).
\end{remark}

\begin{remark}[Choice of $\Stenplus$ vs.\ $S_{10}^-$]
Under $\Spin{10}\supset\Spin{9}$, both Weyl representations restrict
to the same real representation $S_9$ (since $S_9$ has no chirality).
The assignment $\Stenplus$ (rather than $S_{10}^-$) follows
the Boyle convention~\cite{Boyle2020}.  The chirality distinction
will be resolved in Section~\ref{sec:chirality} via the $\Cl{6}$
construction.
\end{remark}

\noindent\textbf{Symmetry upgrade.}
Because $\Vhalf^{\,\mathbb{C}} = \Stenplus$ carries the unique
$\Spin{10}$-module structure extending the $\Spin{9}$ action
(Remark~\ref{rem:why-spin10}), the symmetry of the observer's
description of the Peirce space upgrades:
\begin{equation}\label{eq:spin-upgrade}
  \Spin{9} \;\longrightarrow\; \Spin{10}.
\end{equation}

At the algebra level, the complexification extends from $\Vhalf$ to
the full algebra $\hthree$.  The Peirce multiplication rules
$\Vhalf\circ\Vhalf \subset \Vone\oplus\Vzero$ mean that Jordan
products of elements in $\Vhalf^{\,\mathbb{C}}$ land in
$\Vone^{\,\mathbb{C}}\oplus\Vzero^{\,\mathbb{C}}$: the observer's
complex description of the interface forces a complex description of
every sector reached by interface products.
Concretely, Theorem~\ref{thm:complexification} produces genuinely
complex elements $v = x + iy \in \Vhalf^{\,\mathbb{C}}$ with
$y \neq 0$; their Jordan products $v\circ v$ have imaginary part
$2(x\circ y)_{\Vone\oplus\Vzero}$, which is generically nonzero (it
vanishes only on a measure-zero set of pairs), so the complex
extension of $\Vzero$ and $\Vone$ is used nontrivially, not merely
formally.  Since the Peirce
multiplications generate all of $\hthree$, the observer's effective
algebra is $\hthreeC = \hthree\otimes_\mathbb{R}\mathbb{C}$, not
merely $\Vhalf^{\,\mathbb{C}}$.  The structure group of the
complexified algebra
is~\cite{Baez2002,Yokota2009,Boyle2020}
\begin{equation}\label{eq:F4-E6}
  \Ffour \;\longrightarrow\; \Esix,
  \qquad
  \dim\,\Ffour = 52 \;\to\; \dim\,\Esix = 78.
\end{equation}
These are different kinds of symmetry groups:
$\Ffour = \mathrm{Aut}(\hthree)$ is the automorphism group (preserving
both the Jordan product and the unit element), while
$\Esix = \mathrm{Str}_0(\hthreeC)$ is the structure group of the
complexified algebra (preserving the cubic determinant form up to
scale).  The stabilizer of $E_{11}$ in $\Esix$ is
\begin{equation}\label{eq:stab-E6}
  \mathrm{Stab}_{\Esix}(E_{11}) \;=\; \Spin{10}\times\mathrm{U}(1),
  \qquad
  \dim = 45 + 1 = 46.
\end{equation}

Under $\Spin{10}$, the 27-dimensional representation decomposes as
\begin{equation}\label{eq:27-decomp}
  \mathbf{27}
  \;\to\;
  \underbrace{\mathbf{1}}_{\Vone^{\,\mathbb{C}}}
  \;\oplus\;
  \underbrace{\mathbf{10}}_{\Vzero^{\,\mathbb{C}}}
  \;\oplus\;
  \underbrace{\mathbf{16}}_{\Vhalf^{\,\mathbb{C}} = \Stenplus}\,,
\end{equation}
where $\Vone^{\,\mathbb{C}} \cong \mathbb{C}$ (trivial),
$\Vzero^{\,\mathbb{C}} \cong \mathbf{10}$ (vector of $\Spin{10}$,
verified by the branching $\mathbf{10}|_{\Spin{9}} = \mathbf{9}\oplus\mathbf{1}$),
and $\Vhalf^{\,\mathbb{C}} = \Stenplus$ (Weyl spinor).

\noindent\textbf{Dimension check:}
$1 + 10 + 16 = 27$.\ \checkmark

\begin{remark}[Cross-check with Boyle]
Boyle~\cite{Boyle2020} arrives at the same decomposition
$\mathbf{27} = \mathbf{1}\oplus\mathbf{10}\oplus\mathbf{16}$ under
$\Spin{10}\subset\Esix$ by analyzing the complexified algebra
directly.  The results are consistent; our contribution is to trace the complexification step to the
\emph{combination} of C*-observer nature and the negative-spectrum
Peirce operators (Theorem~\ref{thm:complexification}), rather than
taking it as a mathematical starting point.
\end{remark}

\begin{remark}[Scope of the complexification proof]
\label{rem:complexification-scope}
Theorem~\ref{thm:complexification} establishes a precise algebraic
fact: the CFC computation $\sqrt{T_a}\,T_b\,\sqrt{T_a}$ on anticommuting
Peirce operators produces imaginary coefficients, extending the generated algebra
from $M_{16}(\mathbb{R})$ to $M_{16}(\mathbb{C})$.  This is not a
physical heuristic but a computation within the C*-algebra that the
observer inhabits (Paper~5).  The key input is that the Peirce
operators $T_a$ have eigenvalue $-\tfrac{1}{2} < 0$: in a real
algebra, $\sqrt{-\tfrac{1}{2}}$ is undefined; in the observer's
complex C*-algebra, it is $i/\sqrt{2}$.  The exit from
$\mathbb{R}$ to $\mathbb{C}$ is produced by the \emph{combination}
of the observer's C*-nature (Paper~5) and the spectral properties
of the Peirce operators (the $\hthree$ structure).  Neither input
alone suffices: a real-algebraic observer would have no complex
functional calculus; an operator with non-negative spectrum would
have a real square root.

The upgrade from $\Spin{9}$ to $\Spin{10}$ then follows from the
multiplicity-free branching $\Stenplus|_{\Spin{9}} \cong
S_9^{\,\mathbb{C}}$ (Remark~\ref{rem:why-spin10}): the complexified
space $S_9^{\,\mathbb{C}}$ admits a unique $\Spin{10}$-module
structure extending the $\Spin{9}$ action, whose existence is
guaranteed because $S_9$ is the restriction of $\Stenplus$ to
$\Spin{9} = \mathrm{Stab}_{\Spin{10}}(e_{10})$.

\emph{Interpretive independence.}\enspace
The complexification result requires exactly two mathematical
inputs: (i)~the Peirce operators $T_a$ are faithfully represented
as self-adjoint elements of a C*-algebra $M_n(\mathbb{C})$, and
(ii)~the square root is computed via the C*-algebra continuous
functional calculus.  Input~(i) follows from any faithful
embedding of $M_{16}(\mathbb{R})$ into a complex matrix
algebra.  Input~(ii) is a standard C*-algebra property.  The
physical interpretation of the sequential product (as measurement
composition, per Paper~5) motivates why $\sqrt{T_a}\,T_b\,\sqrt{T_a}$
is the relevant operation, but the algebraic fact that it produces
imaginary coefficients depends only on (i) and (ii).
\end{remark}
